% Cellule d'électrolyse de l'eau à membrane échangeuse de protons (sujet type bac 2, exercice 1)
\begin{tikzpicture}[x=1cm,y=1cm]
\pic at (0,0) {becher={largeur=5.0,hauteur=2.8,niveau=2.35,liquide=pcEau}};
\fill[otMuted!35] (-0.06,0.12) rectangle (0.06,2.3);
\node[pcLegende,anchor=north] at (0,-0.08) {membrane};
\filldraw[pcAppareil] (-1.62,0.35) rectangle (-1.38,2.0);
\filldraw[pcAppareil] (1.38,0.35) rectangle (1.62,2.0);
\foreach \y in {0.7,1.15,1.6}{\draw[otNavy!60,line width=0.4pt] (-1.18,\y) circle (0.07) (-1.08,\y+0.2) circle (0.05);
  \draw[otNavy!60,line width=0.4pt] (1.18,\y) circle (0.07) (1.08,\y+0.2) circle (0.05) (1.25,\y+0.12) circle (0.05);}
\draw[pcFil] (-1.5,2.0) -- (-1.5,3.6) -- (-0.45,3.6);
\draw[pcFil] (1.5,2.0) -- (1.5,3.6) -- (0.45,3.6);
\pcAppareilRond{(0,3.6)}{G}
\node[font=\otSemi\fontsize{9}{10}\selectfont,text=otInk] at (-0.62,3.92) {$+$};
\node[font=\otSemi\fontsize{9}{10}\selectfont,text=otInk] at (0.62,3.92) {$-$};
\draw[otMuted,line width=0.4pt] (-1.62,1.2) -- (-2.75,1.2) (1.62,1.2) -- (2.75,1.2);
\node[pcLegende,anchor=east] at (-2.8,1.2) {électrode A};
\node[pcLegende,anchor=west] at (2.8,1.2) {électrode B};
\draw[otNavy,line width=0.5pt,-{Stealth[length=1.6mm]}] (-1.0,2.45) -- (-1.0,2.95) node[pcLegende,anchor=south,text=otNavy,inner sep=1pt] {gaz A};
\draw[otNavy,line width=0.5pt,-{Stealth[length=1.6mm]}] (1.0,2.45) -- (1.0,2.95) node[pcLegende,anchor=south,text=otNavy,inner sep=1pt] {gaz B};
\node[pcLegende,anchor=west] at (2.8,0.45) {eau en milieu acide};
\end{tikzpicture}
\par\smallskip\legende{Schéma de principe d'une cellule de l'électrolyseur : la membrane laisse passer les ions $\ce{H+}$ d'un compartiment à l'autre.}
